\chapter{119}

\section{Nervi, Genua (IT), im
November}



Sie waren soweit. Sogar die Wohnung hatten sie in Ordnung gebracht.

Einzig ihr letztes Gespräch hing offen in der Luft.

Louis zog den Reissverschluss seines Rucksacks nach unten.

Joh stand unentschlossen im Türrahmen. Er schien mit sich zu ringen.
Sollte er Louis wirklich antworten? Louis konnte nicht nachvollziehen,
was denn so schwierig war. Er liess nicht locker.

«Ich weiss doch, dass du dazu eine Meinung hast, Joh, und dass es deine
persönliche ist.»

«Du willst wirklich, dass ich dir sage, was ich in dieser Situation und
jetzt gerade tun würde?»

«Ja, ich bitte dich darum.»

Seine Stimme hatte unterdessen einen flehenden Ton. Seit Louis
Valentinas Brief gelesen hatte, suchte er ununterbrochen und verzweifelt
nach einer passenden Reaktion.

\sectionbreak

{\letterfont\footnotesize
\indentedblock{
Hi Louis

Wenn du das liest, hat es geklappt. Meine Kusine hat mir versprochen,
dir diesen Brief zu bringen.

Ich konnte leider nicht kommen, weil ich mich schlecht fühle. Ich hoffe,
du bist mir nicht böse...

Nach unserm letzten Anruf hatte ich plötzlich Angst, dir zu begegnen.
Obwohl ich mich seit Tagen auf unser Wiedersehen gefreut habe.

Ich wollte es dir heute ja persönlich sagen, Louis!!! Aber eben...Ich tu
es jetzt mit diesem Brief:

Meine Gefühle für dich sind mehr als freundschaftlich. Ich habe dich
sehr lieb. Du bist so mutig, ehrlich und einfach ein toller Mann.

Nun weisst du es.

Ich hoffe für dich, dass sich die Geschichte mit Giorgia bald ganz
geklärt hat.

Ich denke an dich.

Mi piaci\footnote{\letterfont\scriptsize Ich habe dich gern.},

Valentina
}
}

\sectionbreak

Joh bewegte sich auf Louis zu.

«Pass auf, Louis. Also, du weisst, ich halte wenig bis nichts von
Ratschlägen.»

Louis’ rasches Nicken liess keine Zweifel offen. Wenn er jetzt nicht
gleich eine Antwort bekam, würde er durchdrehen.

Joh kratzte sich am Hals. Dann schaute er Louis direkt an.

«Nichts.»

Louis riss die Augen auf und verschränkte die Arme.

«Nichts? Du würdest nichts tun?»

«Genau.»



Jeder zurückgelegte Kilometer liess Louis’ Nervenkostüm noch enger
werden. Er holte Atem, schaltete runter, beschleunigte den Wagen und
überholte einen breiten, vor sich hin röchelnden Lastwagen. Zurück auf
der rechten Fahrspur blickte er zu Joh hinüber.

«Bald sind wir da, noch fünfzehn Kilometer bis Genua.»

«Jepp.»

Joh verfolgte den Weg auf dem Handy.

«Laut Google-Map sind es von Genua bis Nervi noch knappe zehn Kilometer,
meinen die da.»

Nach einem konzentrierten Blick in den Rückspiegel wechselte Louis die
Spur und überholte eine Reihe Kleinwagen.

\emph{«Nervi} soll Charme haben, schwärmte Giorgias Mutter.
Und sei im November deutlich wärmer als Mailand; zwei bis drei Grad
Unterschied sind nicht ohne.»

Joh streckte sich, so gut es seinem angegurteten Körper möglich war. Er
überschlug die Beine und drehte sich halb zu Louis um.

«Sag mal, wie kam es eigentlich, dass du die Details eures Treffens mit
Giorgias \emph{Mutter} vereinbart hast?»

Louis’ Antwort liess auf sich warten.

«Ach, genau weiss ich es nicht,» Louis drosselte die Geschwindigkeit,
«Giorgia wollte nicht telefonieren, wir schrieben uns. Sie meinte, es
sei einfacher, wenn ich mich mit ihrer Mutter abspreche.»

«Ach ja?»

«Mhm, ich kann nur Vermutungen anstellen. Vielleicht fühlt sie sich zu
schwach oder bekommt beim Sprechen Atemnot, oder…»

«…oder sie ist genau so unsicher wie du.»

«Meinst du?»

Louis wandte sich wieder entschlossen dem Verkehr zu. Der hatte sich
verdichtet. Joh führte ihn mit knappen Anweisungen durch die
Blechlawine.

«Also, der \emph{«Wachturm»} liegt direkt auf der Promenade
\emph{Passeggiata Anita Garibaldi.} Laut Beschreibung führt dieser lange
Weg der Küste entlang,» Joh schob die runtergerutschte Brille hoch, «hat
sie dir genauere Angaben gemacht?»

«Nein, das sei nicht nötig, meinte sie, ich könne ihn nicht verfehlen.»

«Ja, hier, habe ihn.»

Joh scrollte sich durch eine Reihe von Aufnahmen.

Louis wurde still. Sein Herz pochte. Entgegen der angenehmen
Wärme kam eine kalte Welle über ihn. Seine Stimme klang
belegt.

«Ich, ich hoffe, der Weg ist rollstuhlgängig.»
